% Moved from Section 5 to the appendix, Oct 25 2026 (his word: the full 2023 fit goes to the appendix).
\section{The 2023 economy}
\label{sec:quant-2023}

The transition matches only period fertility, so the 2023 economy is an
out-of-sample test of what the decline in $\psi_t$ alone does to the rest of
the economy. Table~\ref{tab:quant-2023} compares the change from 2007 to 2023
in the data and in the model. The model reproduces the shift of first births
to later ages in direction but not in size: the share of women who are
mothers by 25 falls by 13 percentage points in the model and by 18 in the
data, the mean age at first birth rises by 0.9 years against 2.1, and the
share of first births at 30 or later rises by 5 points against 13. Three
changes have the wrong sign. Childlessness among women aged 40 to 44 rises by
8.6 points in the model, while it fell by one point in the data. Ownership
among heads aged 30 to 55 rises by 4.1 points, while it fell by 8.9. The
house price falls by 3.6\%, while the real price rose by 15.6\%. Wealth
relative to earnings, which rose from 4.46 to 7.05 in the PSID, barely moves
in the model. Mean rooms are close to the data, 5.51 against 5.58. The extra
space of families with three or more children falls by 0.03 rooms in both,
but stays too large in the model, 0.54 rooms against 0.36.

\begin{table}[t]
\centering
\caption{The 2023 economy: data and model}
\label{tab:quant-2023}
% [MODEL] experiments/housing_supply/20261021_free_entry/moments_2023_fi/moments_2023_fi.csv,
% free_entry_eta0.63 column (replay of FIN063fi_fw); mothers by 25 = 1 - ever_born_shares_age25['0'].
\small
\begin{tabularx}{\textwidth}{@{}X rr @{\hspace{1.5em}} rr@{}}
\toprule
 & \multicolumn{2}{c}{Data} & \multicolumn{2}{c}{Model} \\
\cmidrule(lr){2-3}\cmidrule(l){4-5}
Moment & 2007 & 2023 & 2007 & 2023 \\
\midrule
Completed fertility, ages 40--44 & -- & 1.92 & 2.10 & 1.69 \\
Childless, ages 40--44 (\%) & 19.8 & 18.8 & 20.1 & 28.7 \\
Exactly one child among mothers, 40--44 (\%) & 21.4 & 22.2 & 21.7 & 22.2 \\
Mean age at first birth (years) & 25.98 & 28.09 & 25.97 & 26.88 \\
Mothers by age 25 (\%) & 45.7 & 27.4 & 45.5 & 32.2 \\
First births at age 30+ (\%) & 24.9 & 37.7 & 24.2 & 28.9 \\
\addlinespace
House price (2007 = 100) & 100 & 115.6 & 100 & 96.4 \\
Mean occupied rooms (capped at 9) & -- & 5.58 & -- & 5.51 \\
Ownership, heads 30--55 (\%) & 67.6 & 58.7 & 66.9 & 71.0 \\
First-birth room response$^{\ddagger}$ & 1.440 & 1.476 & 1.475 & 0.55 \\
Rooms: 3+ versus 1--2 resident children & 0.385 & 0.355 & 0.574 & 0.543 \\
\addlinespace
Wealth / annual gross earnings & 4.46 & 7.05 & 4.35 & 4.38 \\
Annual bequests / wealth (\%) & 0.73 & 0.88 & 0.73 & 0.83 \\
Wealth/income p90 / median, ages 76--84 & 3.52 & 3.45 & 3.27 & 3.29 \\
\bottomrule
\end{tabularx}
\par\smallskip
\begin{minipage}{0.95\linewidth}
\footnotesize\textit{Notes:} Model, 2007: the calibrated stationary
equilibrium (Tables~\ref{tab:quant-moments} and~\ref{tab:quant-untargeted});
model, 2023: the period dated 2023 of the transition in
Section~\ref{sec:quant-transition}, evaluated on the full 2023 distribution.
Data, 2007: the sources of Tables~\ref{tab:quant-moments}
and~\ref{tab:quant-untargeted}; data, 2023: the same sources in their most
recent waves, with wealth over earnings from the PSID 2021 and 2023. House
price: S\&P CoreLogic Case--Shiller national index deflated by the CPI. A
dash marks a 2007 value not measured on the 2023 definition. $^{\ddagger}$ In
2023 the model value compares otherwise identical households with and
without a first birth in the same period, not the event study of the data.
\end{minipage}
\end{table}
\problem{The model and data measures of the first-birth room response have
not been reconciled for the transition: 0.55 is the matched-state comparison,
1.476 the event study. Three 2023 data cells are not yet traced to a 2023
source: bequests 0.88 (recorded only as an external benchmark), old-age
dispersion 3.45 (recorded as PSID 1984--2019) and mothers by 25, 27.4.}
% Oct 25 checks: model 2023 column = moments_2023_fi.csv free_entry_eta0.63 (all rows) and
% validation_2023_model_fi_eta0.63.json (mothers by 25 = 1 - 0.67757). Model 2007 column = calibration
% point center.json / lifecycle_2007_newacct (the transition's own 2007 state, psi 0.177467 vs 0.177405,
% price 0.770687 vs 0.770691, is the same to rounding). Price 2023: FIN063fi_fw 0.742618/0.770686 = 0.964.

The fertility moments reveal the limitations of explaining the decline with
changes in the preference for children alone. In the data, the decline
in period fertility since 2007 came with almost unchanged childlessness among
women aged 40 to 44 (19.8\% in 2004--2006, 18.8\% in 2024) and
with later first births: the mean age at first birth rose from 26.0 to 28.1
and the share of first births at 30 or later from 25 to 38\%. In the
model, a lower $\psi_t$ works on every cohort alive, so 29\% of women
aged 40 to 44 in 2023 are childless and their completed fertility is 1.69,
against 1.92 in the data. The data
show a prominent role for later first births among younger cohorts; the
model instead reduces fertility across the cohorts exposed to the common
preference change. This discrepancy leaves the distribution of preference changes across
cohorts unresolved.

With only preference shocks,
falling fertility lowers the demand for family space, so the model has to
deliver a falling price: 3.6\% below 2007 in 2023. In the data, the real
house price rose by 15.6\% between 2007 and 2023, and real rents by 15
to 19\%. Part of the decline in fertility since 2007 therefore came with
a rising cost of space that the estimation assigns to preferences. The
transition says what housing would have done in response to the fertility
decline, not that housing caused it.
% [DATA] Case-Shiller CSUSHPINSA / CPIAUCSL, 2023/2007 = 1.156
% (historical_supply_comparison.csv); CPI rent of primary residence, FRED.

Figure~\ref{fig:quant-lifecycle-2023} compares life-cycle profiles in 2023.
The model reproduces the rise of ownership with age and the hump in housing
size, but its households buy too early: 43\% of heads aged 18
to 21 and 52\% of those aged 22 to 25 own their home in the
model, against 11\% and 15\% in the ACS. This miss matters
for the credit experiments of Section~\ref{sec:policy}, because it places too many young households
near the margin of buying. First births are also too early:
29\% of first births in the model are to women aged 18 to 21,
against 17\% in the 2021--2023 natality data, and 14\% are at
30 to 33, against 22\%.
% [MODEL] free_entry/moments_2023_fi/lifecycle_by_age_0.63.csv: own 18-21 42.6, 22-25
% 52.2; validation_2023_model_fi_eta0.63.json birth_flow_first shares 18-21 0.288,
% 30-33 0.139. [DATA] actual2023_age_housing_levels.csv: 10.6, 15.3.

\begin{figure}[tbp]
\centering
\includegraphics[width=\linewidth]{../../output/model/calibration/2026-10-24_lifecycle_2023_fi/lifecycle_2023_fi_eta063.pdf}
\caption{Life-cycle profiles in 2023, model and data}
\label{fig:quant-lifecycle-2023}
\par\smallskip
\begin{minipage}{0.92\linewidth}
\footnotesize\textit{Notes:} Homeownership and rooms by age of the household
head, and first births by age of the mother. Model: the 2023 period of the
transition in Section~\ref{sec:quant-transition}. Data: ACS 2023; NCHS first births 2021--2023.
\end{minipage}
\end{figure}
